\chapter{Introduction} \label{sec:introduction}

\section{Motivation}

Many of today’s scientific and machine learning models---including billion-dollar recommendation models and large-scale scientific solvers and simulators---are inherently sparse~\cite{dlrm2020hpca,williams2007spmv}. Yet modern datacenters and supercomputers are optimized for dense computation. This fundamental mismatch calls for architectures that can execute sparse workloads efficiently.

A tensor---a multidimensional array---is said to be \emph{sparse} when most of its entries are zero~\cite{taco2017oopsla}. Many real-world datasets naturally exhibit that structure~\cite{suitesparse2011acm}. For example, the \textit{follow} relation in a social network can be represented as a matrix $A$ where $A_{i,j} = 1$ if user $i$ follows user $j$, and zero otherwise~\cite{naumov2019dlrm}. Since users typically follow only a small fraction of all other users, this matrix is overwhelmingly sparse~\cite{recnmp2020isca}. Similarly, the internal states of scientific solvers and simulations often contain mostly zero-valued elements~\cite{williams2007spmv,nagasaka2018spgemm}.

Because zero elements neither need to be stored nor computed, sparse tensors are represented using compressed formats that track only nonzero values and their coordinates~\cite{barret94csr}. Operating on these formats, however, requires irregular traversals and data-dependent control flow~\cite{spchar2024jpdc}. For example, summation requires merging nonzeros based on their coordinates, which entails frequent data-dependent branch operations~\cite{looplets2023cgo}. Multiplication often requires locating tensor fibers that share a common index, leading to both irregular memory accesses and additional data-dependent control flow.

Traditional datacenters and supercomputers are not designed for such irregular computation. Instead, they are heavily optimized for dense and regular workloads---such as dense matrix multiplication and convolution---where data is laid out contiguously and accessed with predictable patterns~\cite{fpga2022ccpe}. On these architectures, data-dependent branches cause frequent misspeculation and costly pipeline flushes~\cite{spchar2024jpdc}, while irregular memory accesses degrade cache utilization and increase memory latency, ultimately clogging the memory subsystem~\cite{spchar2024jpdc}.

As sparse workloads become increasingly pervasive, there is growing interest in hardware and software solutions tailored to their unique characteristics~\cite{dlrm2020hpca}. Programmable prefetchers represent an initial attempt to adapt traditional general-purpose architectures to irregular access patterns~\cite{yu2015}. Although they require minimal hardware changes, they can only accelerate irregular memory accesses, not control flow, and provide limited speedups~\cite{tmu2023micro}. At the other extreme, specialized accelerators for fixed sparse kernels can deliver substantial performance gains~\cite{zhang2021gamma}, but at the cost of limited flexibility. Reconfigurable architectures broaden the range of supported kernels~\cite{rucker2021capstan}, yet still suffer from classic accelerator drawbacks, including kernel launch overheads and constrained memory capacity.

Decoupled Access--Execute (DAE) architectures~\cite{dae1982isca} are emerging as a promising middle ground that combines the flexibility of general-purpose processors with the efficiency of specialization. DAE architectures separate memory access from computation into distinct units, typically connected through queues. The access unit is optimized for indirect and irregular memory operations, while the execute unit---often a general-purpose core---focuses on computation. Prior DAE proposals can substantially accelerate indirect memory accesses---such as those in sparse matrix--vector multiplication~\cite{spzip2021isca,maple2022isca,planar2021acm}---but they primarily target simple access patterns and do not fundamentally address irregular control flow. As a result, they are not suitable for accelerating general sparse tensor algebra. Therefore, be believe more general yet efficient architectures are needed.

\section{Objectives}

The primary goal of this thesis is to overcome the fundamental challenges of general sparse tensor algebra and enable its efficient execution at scale. To achieve this goal, we propose a hardware--software co-design based on a DAE architecture that spans both processor microarchitecture and compiler stack. Conversely from other solutions in the literature, we aim at designing a system that can efficiently accelerate all sparse and dense tensor algebra operations. We identify the following specific research objectives:

\begin{itemize}
    \item \textbf{Understand the microarchitectural implications of sparsity:} To design efficient hardware, we must first quantify how irregular control flow and memory access impact performance of sparse tensor operations in scientific and machine learning workloads. Although it is well known that sparse kernels such as sparse matrix-vector multiplication perform poorly on traditional architectures~\cite{lin2010spmv,bell2009spmv,sadi2019spmv,fowers2014spmv,alappat2022spmv,grigoras2015spmv,umuroglu2014spmv,williams2007spmv,nurvitadhi2015spmv}, our objective is to generalize this understanding to a broader class of sparse tensor algebra operations, each with distinct characteristics and performance challenges, and to delve deeper into the fundamental causes of their inefficiency. By characterizing how these workloads stress modern processors, we establish a detailed understanding of why existing architectures fall short and what requirements future designs must satisfy to scale effectively.

    \item \textbf{Design an efficient DAE architecture for sparse tensor algebra:} Building on this analysis, we aim to develop a DAE architecture capable of overcoming all the identified bottlenecks. Our objective is to design a general and scalable solution that accelerates all forms of sparse tensor algebra while outperforming traditional CPUs and GPUs in both raw performance and performance per watt.

    \item \textbf{Provide seamless integration into state-of-the-art machine learning frameworks:} To fully unlock the benefits of the proposed architecture, we aim to design a compiler stack that integrates naturally with existing tensor programming ecosystems such as PyTorch and TensorFlow, enabling the automatic generation of optimized DAE code without imposing any additional programmability burden on users.
\end{itemize}

\section{Contributions}

This thesis addresses these research objectives through the following contributions.

\subsection{Performance Analysis of Sparse Workloads}

The first contribution of this thesis is a comprehensive analysis of sparse tensor algebra operations across scientific and machine learning workloads. In \Cref{sec:architectural-implications-of-sparsity}, we show that irregular tensor traversal and merging are fundamental bottlenecks in traditional architectures. We first benchmark off-the-shelf CPUs and GPUs, and demonstrate that irregular memory accesses and control flow can stall computation of several sparse tensor algebra workloads for up to 99\% of the time. Then, we characterize representative algorithms and inputs, demonstrating how characteristics such as spatio-temporal locality relate to these bottlenecks. Finally, we model scaling of performance, power, and area to demonstrate that scaling up and out traditional compute cores is not a viable path toward efficient sparse tensor algebra.
This contribution led to the following publications:
\begin{itemize}
\item Francesco Sgherzi, \textbf{Marco Siracusa}, Ivan Fernandez, Adrià Armejach, Miquel Moretó. 2024. \textit{SpChar: Characterizing the Sparse Puzzle via Decision Trees}. \textit{Journal of Parallel and Distributed Computing}~\cite{spchar2024jpdc}.
\item \textbf{Marco Siracusa}, Miquel Moretó, Adrià Armejach. Accelerating Graph Machine Learning with Decoupled Memory-Access--Execute Architectures. \textit{Under submission}.
\end{itemize}

\subsection{The Tensor Marshaling Unit}

The second contribution is the \emph{Tensor Marshaling Unit} (TMU), a specialized hardware unit that offloads and decouples tensor traversal and merging from the compute core, enabling scalable execution. Unlike prior specialized designs, the TMU (1)~accelerates \emph{all} sparse and dense tensor algebra expressions, and (2)~integrates seamlessly with modern vector units to exploit structured and unstructured parallelism. In \Cref{sec:the-tmu}, we present the TMU architecture, its integration into a multicore DAE processor composed of TMUs and CPUs, and a detailed evaluation of performance, power, and area. The resulting DAE processor outperforms traditional CPUs by up to an order of magnitude and GPUs by up to 6$\times$ in both performance and performance per watt.

This contribution led to the following publication:
\begin{itemize}
    \item \textbf{Marco Siracusa}, Víctor Soria-Pardos, Francesco Sgherzi, Joshua Randall, Douglas J. Joseph, Miquel Moretó Planas, Adrià Armejach. 2023. \textit{A Tensor Marshaling Unit for Sparse Tensor Algebra on General-Purpose Processors}. In Proceedings of the \textit{56th Annual IEEE/ACM International Symposium on Microarchitecture} (MICRO '23)~\cite{tmu2023micro}.
\end{itemize}
and patent:
\begin{itemize}
    \item \textbf{Marco Siracusa}, Joshua Randall, Douglas James Joseph, Miquel Moretó Planas, Adrià Armejach Sanosa. March 6, 2025. \textit{Data structure marshalling unit}. US 20250077499A1. Barcelona Supercomputing Center; Arm Ltd.~\cite{siracusa2025dsmu}.
\end{itemize}

\subsection{The Ember Compiler}

The third contribution is the \emph{Ember compiler}, which lowers PyTorch and TensorFlow programs to optimized TMU--CPU code. In \Cref{sec:dae-compilation}, we first show that fully decoupled access and execution introduce significant programmability challenges. We then present Ember’s intermediate representations and compilation strategies, specifically designed for DAE optimization and tightly integrated with modern tensor compilers. Finally, we demonstrate that Ember achieves performance comparable to hand-optimized implementations, fully unlocking the TMU’s potential at no programmability cost.

This contribution led to the following publication:
\begin{itemize}
    \item \textbf{Marco Siracusa}, Olivia Hsu, Víctor Soria-Pardos, Joshua Randall, Arnaud Grasset, Eric Biscondi, Doug Joseph, Randy Allen, Fredrik Kjolstad, Miquel Moretó Planas, Adrià Armejach. 2026. \textit{Ember: A Compiler for Embedding Operations on Decoupled Access--Execute Architectures}. In Proceedings of the \textit{24th Annual ACM/IEEE International Symposium on Code Generation and Optimization} (CGO '26)~\cite{ember2026cgo}.
\end{itemize}

\section{Thesis Structure}

The remainder of this thesis is organized as follows:
\begin{itemize}
\item \Cref{sec:background-and-motivation} provides background on sparse tensor algebra and its role in scientific and machine learning workloads.
\item \Cref{sec:related_work} reviews prior hardware and software approaches for sparse tensor algebra and motivates the need for new DAE systems such as the TMU--CPU architecture proposed in this thesis.
\item \Cref{sec:architectural-implications-of-sparsity} presents the first contribution: a quantitative characterization of sparse workloads and the architectural insights that guide our co-design.
\item \Cref{sec:the-tmu} presents the Tensor Marshaling Unit, its integration into a DAE multicore processor, and a comprehensive evaluation.
\item \Cref{sec:dae-compilation} presents the Ember compiler, the challenges of decoupled execution, and the corresponding compilation techniques and evaluation.
\item \Cref{sec:conclusions} summarizes the key findings of this thesis.
\item \Cref{sec:future-work} outlines promising directions for future research.
\end{itemize}